\documentclass[12pt]{article}
\usepackage[spanish]{babel}
\usepackage[utf8]{inputenc}
\usepackage[T1]{fontenc}
\usepackage{amsmath,amssymb}
\usepackage[margin=2cm]{geometry}

\begin{document}

\begin{enumerate}
    \item \textbf{(30\%)} Muestre la siguiente identidad (puede usar coordenadas cartesianas):
    \[
    \nabla^2 \vec{A} = \nabla(\nabla \cdot \vec{A}) - \nabla \times (\nabla \times \vec{A})
    \]

    \item Las ecuaciones para el campo electromagnético pueden escribirse en términos de potenciales, el potencial escalar $\phi$ y el potencial vectorial $\vec{A}$. Conociendo estos potenciales y una condición adicional (llamada condición de Gauge), podemos conocer el campo electromagnético.

    \begin{enumerate}
        \item[(a)] \textbf{(20\%)} Use la notación de Levi-Civita y pruebe las siguientes identidades vectoriales:
        \begin{itemize}
            \item
            \[
            \vec{D} = -\nabla \phi \Rightarrow \nabla \times \vec{D} = 0
            \]
            \item
            \[
            \vec{D} = \nabla \times \vec{A} \Rightarrow \nabla \cdot \vec{D} = 0
            \]
            \item
            \[
            \nabla \times (\nabla \times \vec{D}) - \nabla(\nabla \cdot \vec{D}) + \nabla^2 \vec{D} = 0
            \]
        \end{itemize}

        \item[(b)] \textbf{(20\%)} Utilice el resultado anterior y las ecuaciones de Maxwell,
        \begin{align}
        \nabla \cdot \vec{E}(\vec{r},t) &= \frac{\rho(\vec{r},t)}{\varepsilon_0}, \\
        \nabla \cdot \vec{B}(\vec{r},t) &= 0, \\
        \nabla \times \vec{E}(\vec{r},t) &= -\frac{\partial \vec{B}(\vec{r},t)}{\partial t}, \\
        \nabla \times \vec{B}(\vec{r},t) &= \mu_0 \vec{J}(\vec{r},t) + \mu_0\varepsilon_0 \frac{\partial \vec{E}(\vec{r},t)}{\partial t},
        \end{align}
        para mostrar que los campos $\vec{E}(\vec{r},t)$ y $\vec{B}(\vec{r},t)$ son derivables de un potencial escalar $\phi(\vec{r},t)$ y de un potencial vectorial $\vec{A}(\vec{r},t)$ en la forma
        \[
        \vec{B}(\vec{r},t)=\nabla \times \vec{A}(\vec{r},t),
        \qquad
        \vec{E}(\vec{r},t)=-\nabla\phi(\vec{r},t)-\frac{\partial \vec{A}(\vec{r},t)}{\partial t}.
        \]

        \item[(c)] \textbf{(30\%)} Reemplace las expresiones encontradas para $\vec{A}(\vec{r},t)$ y $\phi(\vec{r},t)$ en las ecuaciones de Maxwell y muestre que los potenciales electromagnéticos satisfacen las siguientes ecuaciones
        \begin{align}
        \nabla^2\phi(\vec{r},t) + \frac{\partial}{\partial t}\big(\nabla \cdot \vec{A}(\vec{r},t)\big)
        &= -\frac{\rho(\vec{r},t)}{\varepsilon_0}, \\
        \nabla\left(\nabla \cdot \vec{A}(\vec{r},t) + \mu_0\varepsilon_0\frac{\partial \phi(\vec{r},t)}{\partial t}\right)
        + \nabla^2\vec{A}(\vec{r},t)
        - \mu_0\varepsilon_0\frac{\partial^2 \vec{A}(\vec{r},t)}{\partial t^2}
        &= -\mu_0\vec{J}(\vec{r},t).
        \end{align}

        ¿Qué condición matemática deben satisfacer los potenciales para que las ecuaciones anteriores se desacoplen? Reescriba de nuevo dichas ecuaciones pero en su forma desacoplada. ¿Qué puede concluir de su resultado?
    \end{enumerate}
\end{enumerate}

\end{document}